\metatitle{K-theory}
\metadescription{A short introduction to K-theory, Grothendieck groups, sheaves, and K-theoretic Schubert calculus.}

\section[kTheory]{K-theory}

\defin{K-theory} studies a space by looking at vector bundles, or more
generally sheaves, on that space.  The basic construction is a
\defin{Grothendieck group}: start with geometric objects, add formal
differences, and impose the relation that a middle object is the sum of its
subobject and quotient.

For algebraic combinatorics, the main reason to care is that K-theory is the
natural home for \hyperref[grothendieck]{Grothendieck polynomials}.  Ordinary
\hyperref[schubertCalculus]{Schubert calculus} multiplies fundamental classes
of \hyperref[schubertVarietyDefinition]{Schubert varieties} in cohomology.
K-theoretic Schubert calculus instead multiplies the classes of their
structure sheaves, and the resulting polynomial representatives are
Grothendieck polynomials rather than
\hyperref[schubert]{Schubert polynomials}.


\section[grothendieckGroup]{The Grothendieck group idea}

Suppose we have a class of objects with exact sequences, for instance vector
bundles on a space $X$.  The Grothendieck group is generated by symbols
$[E]$, one for each vector bundle $E$, with the relation
\[
  [E]=[E']+[E'']
\]
whenever
\[
  0\longrightarrow E'\longrightarrow E\longrightarrow E''
  \longrightarrow 0
\]
is exact.  Direct sum gives addition:
\[
  [E\oplus F]=[E]+[F].
\]
Tensor product gives multiplication:
\[
  [E]\cdot[F]=[E\otimes F].
\]
The resulting ring is often denoted $K^0(X)$.

\begin{example*}[The point]
Every complex vector bundle on a point is just a finite-dimensional complex
vector space.  Hence
\[
  K^0(\mathrm{pt})\cong \setZ,
\]
where a vector space is sent to its dimension.
\end{example*}

This should be compared with cohomology.  Cohomology remembers cycle classes
and intersection numbers.  K-theory remembers bundle and sheaf classes.  There
are maps from K-theory to cohomology, such as the Chern character, but the
multiplication in K-theory keeps information which is invisible if we only
look at fundamental cohomology classes.


\section[kTheoryCoherentSheaves]{Coherent sheaves}

In algebraic geometry one often uses coherent sheaves instead of only vector
bundles.  The group built from coherent sheaves is usually denoted $K_0(X)$.
For a short exact sequence
\[
  0\longrightarrow \mathcal F'\longrightarrow \mathcal F
  \longrightarrow \mathcal F''\longrightarrow 0
\]
we impose
\[
  [\mathcal F]=[\mathcal F']+[\mathcal F''].
\]

\begin{example*}[A divisor]
Let $D\subset X$ be a divisor.  The standard exact sequence
\[
  0\longrightarrow \mathcal O_X(-D)\longrightarrow \mathcal O_X
  \longrightarrow \mathcal O_D\longrightarrow 0
\]
gives
\[
  [\mathcal O_D]=[\mathcal O_X]-[\mathcal O_X(-D)]
  \qquad\text{in } K_0(X).
\]
Thus subvarieties naturally appear through structure sheaves, but the
relations involve subtraction.  This is one source of the signs that appear in
K-theoretic Schubert calculus.
\end{example*}


\section[kTheoreticSchubertCalculus]{K-theoretic Schubert calculus}

Let $X$ be a Grassmannian or \hyperref[flagVarietyDefinition]{flag variety},
and let $X_w$ be a \hyperref[schubertVarietyDefinition]{Schubert variety}.
In cohomology we use the Schubert class $[X_w]$.  In K-theory we use
the class
\[
  [\mathcal O_{X_w}]
\]
of its structure sheaf.  These classes form a basis of the K-theory of the
\hyperref[flagVarietyDefinition]{flag variety} or Grassmannian.  Since these varieties are smooth, the product
of coherent-sheaf classes can be defined using the derived tensor product,
and multiplication has the form
\[
  [\mathcal O_{X_u}]\,[\mathcal O_{X_v}]
  =
  \sum_w c_{u,v}^w[\mathcal O_{X_w}].
\]
The constants $c_{u,v}^w$ are K-theoretic analogues of Schubert structure
constants.

In the Grassmannian case, these constants are computed by
K-theoretic Littlewood--Richardson rules.  A key combinatorial model is given
by set-valued tableaux, introduced by \name{Anders S. Buch} for the K-theory
of Grassmannians \cite{Buch2002}.  On the polynomial side, \name{Alain
Lascoux}'s Grothendieck polynomials represent K-theory Schubert classes,
namely the structure-sheaf classes of Schubert varieties in the Grothendieck
ring of the flag variety \cite{Lascoux1990Grothendieck}.


\section[kTheoryHowToRead]{How to read more}

For a general entry point to topological K-theory, see \name{M. F. Atiyah}'s
book \cite{Atiyah1967KTheory}.  For Schubert calculus and the Grassmannian
side, \name{William Fulton}'s book on Young tableaux \cite{Fulton1997} and
\name{Maria Gillespie}'s introduction \cite{Gillespie2019} are good first
reads.  For the K-theoretic Grassmannian story and set-valued tableaux, see
Buch's paper \cite{Buch2002}.  For the polynomial representatives used
throughout this site, continue with
\hyperref[grothendieck]{Grothendieck polynomials}.
